\subsection{Begin a reporting period with live claims}
\label{sec:complete-claims-opening}

A new reporting period need not begin with an empty receivables register.
Customers may still owe amounts from sales in the preceding year. The closing
procedure carries those claims forward with their identifiers, issue dates,
remaining amounts and assessment history. An opening rule that requires an
empty register would reject this ordinary continuation even when the carried
book and opening statements are internally consistent.

The retained examples make that conflict concrete. Two company charts carry
one live claim with a remaining amount of 20 into 1 January 2028. Two further
programs carry three claims whose remaining amounts are 13 each, giving a
receivables balance of 39. The claims were issued during 2027 and assessed on
31 December 2027. The new reporting period ends exclusively on 1 January 2029;
its period flows are empty, while the claim records and previously seen event
identities remain. The original continuation contains three commands in each
one-claim case and two in each three-claim case. Those ten commands supply an
exact comparison for the opening extension.

The selected empty-opening template rejects all four books at
\code{empty\_claims}. Removing that stage alone exposes a second empty-specific
condition: \code{claim\_total} requires the receivables balance to be zero.
Removing both conditions makes the remaining opening checks pass, but it also
admits five deliberately corrupted books. The corruptions are a wrong
receivables total, a negative remaining claim, a missing sale identity, an
issue date after the opening and an inconsistent overdue amount. The full
selected state invariant rejects each corruption. Thus the required change
is not simply to permit a nonempty map. It is to retain the relationship
between every supplied claim and the book as a whole.

\code{CompleteClaimsOpening/1} makes that change under an independently
selected contract. The selection identifies the original empty-opening
template and the full state invariant. The new contract carries both
definitions, their canonical identities, the exact two excluded stage names
and an explicit description of its scope. Before deriving the revised
structural template, the adapter checks that \code{empty\_claims} is the
reviewed equality with the empty claim map and that \code{claim\_total} is
the reviewed zero-receivables predicate. Every other stage remains unchanged.
A newly supplied predicate does not acquire reviewed status merely because
the caller computes its digest.

The accompanying \code{SuppliedClaimPopulation/1} declaration states whether
the supplied internal claim map is complete, partial or unknown. It names
\code{/claims}, identifies the exact opening book and supplies provenance.
Only an explicit \code{complete} declaration proceeds under this profile;
partial and unknown declarations return unresolved before any command runs.
Completeness here concerns the supplied internal state. A populated map, a
balanced book and a successful continuation cannot establish that a real
issuer has disclosed every claim. That different question still belongs to
source-constrained reconstruction.

Admission first runs the retained structural and opening-report checks on the
same input. It then evaluates the full selected invariant over the unchanged
book, schemas and memory definition. The invariant checks the aggregate
receivables relationship and the claim conditions that the simple relaxation
missed. A malformed schema or incomplete statement mapping therefore still
fails its original check. A mid-period snapshot also remains unsuitable as a
new-period opening even if its claims satisfy the invariant. This profile
deliberately checks structure and the opening report before the full
invariant; it does not promise an older profile's failure precedence when
several conditions fail together.

The program wrapper, \code{AdmittedCompleteClaimsProgram/1}, retains the
complete admission result. After admission succeeds, it delegates to the
existing opening-to-program binding using the derived structural template
and its corresponding internal selection. That binding again checks shared
facts, schemas, mapping, memory definition, roles, policy, rules, invariant,
complete projected commands and event-input order before the unchanged
composer runs. A caller-supplied success flag cannot replace these checks.
The retained delegated result exposes the requests, transitions, reports,
carry state, trace and any stopped suffix.

The same distinction is explicit in company source.
\code{CompleteClaimsCompanyText/1} surrounds the full selected company
source with a reviewed-selection identity, population declaration and contract
residuals. The declaration below supplies the population premise in a checked
carried example; the opening-book identity is computed from the exact imported
book rather than typed a second time.

\begin{quote}\small\ttfamily
population "SuppliedClaimPopulation/1" complete "/claims"\\
\hspace*{1em}"Exact retained synthetic carry; complete supplied internal map";
\end{quote}

This display wraps one source declaration for readability. The complete
program begins with the text and target profiles, then names the expected
reviewed-selection digest. The caller supplies that reviewed selection
separately. Between \code{begin\_selected\_company} and
\code{end\_selected\_company}, the author writes the complete opening,
invariant, eligibility, recipe and command source described in
Section~\ref{sec:selected-company-source}. Its only imported resources remain
the opening book, schemas, statement mapping and memory definition. An extra
resource containing an opaque algorithm is refused.

Lowering checks that both the named identity and the nested source selection
match the independently supplied selection. Editing an opening condition or
invariant cannot silently replace that reviewed choice. Lowering also verifies
the structural contract, but does not admit the opening or run commands.
Execution invokes admission anew. UTF-8 bytes and line endings remain part of
source identity; nested errors retain their original parent-source locations.
The invariant parser supplies declaration locations, while the opening lexer
can identify an offending token. The outer source preserves that difference
instead of asserting precision the inner parser does not provide.

Four carried and four empty-opening source programs reproduce the complete
typed results. The carried programs retain all ten continuation commands;
the empty-opening programs retain all 44 commands. The five corrupted states
still refuse before execution. Native lifecycle callbacks and the operating
method builder are disabled throughout these comparisons. The result
comparison permits differences only in the outer and delegated
\code{program\_order\_sha256} fields: both record the full program's
serialization order, which source reconstruction can change while preserving
canonical content. Both identities are independently recomputed and retained.
Every other field, including complete execution and event-input-order identity,
matches the typed comparison.

The checked successor therefore admits these complete supplied carried books
without weakening the selected invariant. It leaves partial actual histories,
arbitrary alternative financial rules and accounting applicability unresolved.
The carried-state audit, typed contracts and full source comparisons are kept
together with their explicit refusal cases under
\begingroup\raggedright
\path{data/accounting-master/nonempty-opening-audit-20261026/},
\path{data/accounting-master/complete-claims-opening-20261026/}
and \path{data/accounting-master/complete-claims-source-20261026/}.
\par\endgroup
